\documentclass[11pt]{article}
\usepackage[margin=1in]{geometry}
\usepackage{amsmath,amssymb,amsthm,hyperref}
\newtheorem{lemma}{Lemma}
\title{ARC Joint-Simulation Audit: Recovered Barrier and Exact Two-Stage Target}
\author{PNP ARC research run}
\date{October 7, 2026}
\begin{document}
\maketitle
\paragraph{Status.} Source reconciliation plus proved conditional arithmetic.
No SAT-specific closure theorem or P versus NP separation is claimed.

\section{Review before construction}
The visible conversation establishes the two proof spines and the already
proved exponent iteration. A fresh GitHub read of \texttt{main4.tex}, SHA
\texttt{d1274b98fcd8c41cde544171e739d8893bd3fb6b}, confirms that finite
return-cycle contraction has already been added. We do not claim that model
as a new result. The personal-context search service returned unavailable,
so this run cannot claim an exhaustive review of other conversations.

The Drive document SIM-4-0008, dated September 28, was fetched in full. It
already establishes the exact guard that a fast original machine and a
low-space simulator cannot have their separate coordinates combined into
one joint bound. It explicitly retires the shortcut from a fast contractor
to a favorable joint-space contractor via Williams's simulator. This is a
recovered solved logical barrier, not a newly discovered open question.
Its remaining target is a model-matched simulator with both coordinates
certified on the same realization.

\section{Current primary-literature check}
Williams's 2025 result supplies a multitape time-to-space simulation with
space $O(\sqrt{t\log t})$ \cite{williams2025joint}. Its Section 5 asks about
joint time-space tradeoffs and identifies costly low-degree-extension
evaluation. The theorem therefore supplies no time-preserving joint pair
for this ARC application. Buss and Williams delimit the formalized
alternation-trading methods; that result does not rule out all SAT-specific
improvements \cite{busslimitsjoint}. These sources justify retaining the
missing joint implication rather than treating a space-only simulation as
its proof.

\section{Exact two-stage candidate}
This section isolates a useful numerical target, without attributing either
phase to a known general SAT algorithm. Suppose a model-matched joint
simulator, on input achievable time exponent $t$, provides one witness with
\[
 (t_1,s_1)\le(qt,\alpha t),\qquad q\ge1,\quad \alpha\ge0.
\]
Suppose a second, SAT-specific transformation provides another single joint
witness with
\[
 (t_2,s_2)\le(a+b t_1+c s_1,\ d+s_1),
 \qquad a,b,c,d\ge0.
\]
Both displays are hypotheses. In particular, the second coordinate cannot
be copied from the first phase without an implementation or existence proof
that the second phase preserves it. All encoding and simulation overhead is
included in the displayed constants.

Let $\theta=bq+c\alpha$. A complete return has envelope
\[
 t_{r+1}=a+\theta t_r,\qquad s_{r+1}=d+\alpha t_r.
\]
\begin{lemma}
If $\theta<1$, the envelope converges to
\[
 t_*=\frac{a}{1-\theta},\qquad s_*=d+\alpha t_*.
\]
If in addition
\[
 \frac{(1+\alpha)a}{1-bq-c\alpha}+d<2\cos(\pi/7),
 \tag{J}
\]
then finitely many complete returns from any finite starting time exponent
produce an envelope in the product-endpoint forbidden region.
\end{lemma}
\begin{proof}
Induction gives $t_r=t_*+\theta^r(t_0-t_*)$ and
$s_r=d+\alpha[t_*+\theta^{r-1}(t_0-t_*)]$ for $r\ge1$.
Their sum converges to the left side of (J). Strict inequality yields finite
crossing. Componentwise domination transfers the envelope bound to actual
witnesses when both phase hypotheses are established.
\end{proof}
The useful remaining target is now explicit. A simulator's time overhead
$q$ is multiplied by the contractor's slope $b$, and inherited space adds
$c\alpha$ to the time slope. Obtaining $bq+c\alpha<1$ is only the stability
condition; (J) separately controls the endpoint.

Even a hypothetical perfect-time square-root-space simulation alone would
not erase an arbitrary time exponent: $(t,s)\mapsto(t,t/2)$ has eigenvalue
one, and repeated application leaves $t$ unchanged. Its product exponent
stabilizes at $3t/2$, which need not approach the target. This arithmetic
failure is stronger than merely observing that the published simulator's
time cost is unspecified for our needs. A separate time-improvement phase
is necessary for this candidate architecture.

\section{Bounded test and hostile cases}
The complete script \texttt{arc\_joint\_simulation\_test.py} uses rational
arithmetic to test 108 parameter combinations. It verifies the exact
recurrence formulas for stable cases and rejects mismatched witness or
model tags. It also checks that repeated perfect-time square-root-space
simulation cannot reduce an initial time exponent ten. The run completed
1,602 checks in about 0.009 seconds with a 25-second internal deadline.
Results are recorded in \texttt{arc\_joint\_simulation\_results.json}.
Floating endpoint margins are numerical illustrations only; rational fixed
points are retained in the output. No tested case instantiates a general
SAT transformation.

\section{Precise handoff}
Preserve the recovered same-realization guard as established. Search for or
construct phase one with a joint bound rather than a space-only theorem,
and phase two with a proved inherited-space guarantee. An actual
source-dependent transformation meeting (J), or another independently proved
SAT closure, is still required. A nonrelativizing ingredient does not by
itself provide those quantitative implications.

\begin{thebibliography}{9}
\bibitem{williams2025joint} R. Ryan Williams.
\emph{Simulating Time With Square-Root Space}. STOC 2025.
\url{https://arxiv.org/html/2502.17779v1}. Theorem 1.1 and Section 5.
\bibitem{busslimitsjoint} S. Buss and R. Williams.
\emph{Limits on Alternation-Trading Proofs for Time-Space Lower Bounds}.
Author publication page and manuscript.
\url{https://mathweb.ucsd.edu/~sbuss/ResearchWeb/npProofLimits/}.
\bibitem{sim40008} PNP project. SIM-4-0008.
\emph{Square-Root-Space Shadow Fails SAME\_REALIZATION, Exposes Nonrelativizing Tradeoff Target}.
September 28, 2026.
\url{https://docs.google.com/document/d/1LGCd_cAT-nC8s3iyIKqvhkKJAi7gGLYmFUo88r0n16s/edit}.
\bibitem{arcgithub} PNP project. \texttt{main4.tex}, inspected revision
\texttt{d1274b98fcd8c41cde544171e739d8893bd3fb6b}.
\url{https://github.com/j15marti/pnp/blob/main/main4.tex}.
\end{thebibliography}
\end{document}
